%!TEX TS-program = lualatex
\documentclass[aspectratio=169]{beamer}

% Настройка языка и шрифтов для LuaLaTeX
\usepackage{fontspec}
\usepackage{polyglossia}
\setmainlanguage{russian}
\setotherlanguage{english}
\setmainfont{DejaVu Serif}
\setsansfont{DejaVu Sans}
\setmonofont{DejaVu Sans Mono}

% Оформление
\usetheme{Madrid}
\usecolortheme{whale}

% Дополнительные пакеты
\usepackage{graphicx}
\usepackage{listings}
\usepackage{hyperref}

\lstset{
    basicstyle=\ttfamily\small,
    breaklines=true,
    frame=single
}

% Метаданные
\title[Инструменты визуализации]{Разработка научных презентаций и постеров в LaTeX}
\author[Тоц Л. А.]{Тоц Леонид Александрович}
\institute[ИВТ-2]{Группа ИВТ-2}
\date{\today}

\begin{document}

% Титульный слайд
\frame{\titlepage}

% Слайд 1: Списки и оверлеи
\begin{frame}{Основы Beamer: Списки и наложения}
    \begin{itemize}
        \item<1-> Презентации компилируются в стандартный формат PDF.
        \item<2-> Наложения (\texttt{overlays}) позволяют пошагово выводить тезисы.
        \item<3-> Поддержка динамического акцентирования внимания слушателей.
    \end{itemize}
\end{frame}

% Слайд 2: Мультиколонки и графика
\begin{frame}{Структурирование: Колонки и графика}
    \begin{columns}[c]
        \begin{column}{0.48\textwidth}
            \textbf{Преимущества разделения на колонки:}
            \begin{itemize}
                \item Рациональное использование площади экрана 16:9;
                \item Сопоставление кода и графиков;
                \item Улучшенная читаемость данных.
            \end{itemize}
        \end{column}
        \begin{column}{0.48\textwidth}
            \begin{center}
                % Заглушка стандартного изображения LaTeX
                \includegraphics[width=0.85\linewidth]{example-image}
            \end{center}
        \end{column}
    \end{columns}
\end{frame}

% Слайд 3: Блоки, код и ссылки
\begin{frame}[fragile]{Информационные блоки и листинги}
    \begin{block}{Ключевой вывод}
        Инструменты \texttt{listings} и блоки окружения позволяют оформлять листинги без потери форматирования.
    \end{block}

    \vspace{0.5em}
    \begin{lstlisting}[language=Python]
def calculate_metrics(data):
    # Пример обработки выборки
    return sum(data) / len(data)
    \end{lstlisting}

    \vspace{0.5em}
    Материалы и документация доступны по ссылке: \href{https://ctan.org/pkg/beamer}{\color{blue}\underline{CTAN: Beamer Package}}.
\end{frame}

\end{document}